% DRAFT -- new section, not yet included in main.tex.
%
% Requested Sep 6: "a dedicated section on what was gained by having a fully
% agent-driven design loop", covering the cost_aware_tool_score finding, the
% leaderboard-forensics capability, and anything else relevant.
%
% Written without regard to the page limit, per instruction.  Measured, not
% estimated: \input{} before Sect. 5 takes the body from 6 pages to 8, with
% references then spanning p9-p10.  Compiles with 0 errors and 0 unresolved
% references.  See the note at the foot of this file for integration options.
%
% Every number below was recomputed from val_metrics_S2.json and
% best_val_metrics_sep_6.json on Sep 6 and is reproducible from those two
% files; neither is committed (both are gitignored -- see .gitignore).

\section{What the Loop Bought}
\label{sec:gained}

Section~\ref{sec:human} is an inventory of the loop's failures, and read alone it
argues for a human. This section is its counterweight: four things the loop did that
we do not think a conventionally-run entry would have done, not because a human
could not do them, but because each has a cost profile that an agent changes. We
state the asymmetry in each case rather than claiming a capability.

\paragraph{It read the scorer instead of sampling it.}
The released \code{evaluate.py} is roughly 1{,}300 lines. The loop's first act was to
read all of it and derive the case score's six weights, the ranking formula, and the
gate semantics from source --- before running a single experiment. Two of the
design's load-bearing decisions follow directly from that read and from nothing else:
that $0.75$ of the case score is deterministic given the decision, which is what
makes a template competitive with a language model; and that Task~3's ranking score
is the C-index \emph{alone}, so its rationale is unscored. The second is worth
stating concretely, because it is the kind of thing a leaderboard cannot teach you:
our Task~3 rationale scores $0.2739$ against the public leader's $0.6565$, and the
difference is worth exactly zero. An entry that had inferred the metric from
submissions would have spent slots discovering this. The asymmetry is not insight
but \emph{tedium tolerance}: reading 1{,}300 lines of scoring code carefully enough
to trust a claim about component weights is a task humans reliably skim and agents
reliably do not.

That same read produced the \code{cost\_aware\_tool\_score} finding reported in
Sect.~\ref{sec:human}, and the validation
metrics let us confirm it against an independent entrant rather than only by
argument: across all 86 Task~1 and~2 validation cases, both our entry and the
public leader score \code{mean\_tool\_score} of exactly $1.0000$. A component
carrying $0.15$ of the case score --- the third-largest weight --- separates the top
two entries by nothing at all, which is what a precision with $|A| = \emptyset$
mapped to $1.0$ predicts. We reported the property to the organizers rather than
exploiting it, and note here only that it is measurably non-discriminative, not
merely theoretically so.

\paragraph{It could do forensics on a competitor's public metrics.}
The platform publishes per-case component scores for leaderboard entries. When the
top entry stood at $0.8719$ against our $0.8061$, the loop decomposed the $0.0658$
gap in a single pass --- 109 cases, six components, restricted to the cases both
systems' gates passed so the comparison is like-for-like. The result reframed the
remaining work:

\begin{center}
\footnotesize
\begin{tabular}{lccc}
\toprule
& Task 1 & Task 2 & Task 3 \\
\midrule
gap in ranking score        & $.0512$ & $.0802$ & $.0661$ \\
$\times$ task weight        & $.0205$ & $.0321$ & $.0132$ \\
\midrule
\emph{of which} reasoning trace & $.0035$ & $.0107$ & $.0000$ \\
\emph{of which} decision / ordering & $.0170$ & $.0214$ & $.0132$ \\
\bottomrule
\end{tabular}
\end{center}

\noindent
\textbf{Seventy-eight per cent of the gap to the leader is decisions and ordering;
$22\%$ is the reasoning trace} --- although the trace carries $0.75$ of the case
score. This is the paper's thesis measured against another team's system rather than
our own, and it is the strongest form of the claim we can make: the trace is largely
a supervised target that a careful template saturates, and the residual difficulty
is in the decisions.

The per-case view sharpened it further. On Task~1 the leader's four errors are a
strict \emph{subset} of our six; all six are the same failure, a ground truth of
\code{yes} called \code{no}; and on the 44 cases both gates passed, our total
reasoning deficit is $0.0197$ of case score, of which we are actually \emph{ahead} on
section grounding ($0.6959$ against $0.6291$). Two biopsy calls out of fifty separate
the two entries on Task~1. On Task~2 the reasoning gap is real but concentrated:
$0.0565$, of which confidence conditioning is $0.0206$ --- the one component where a
constant, cross-validated optimum on our training cohort is visibly beaten on
validation.

We want to be exact about what this changed, because the tempting claim is the wrong
one. It supplied no method: the S3 change described in Sect.~\ref{sec:method} was
derived from our own training labels, and we deliberately did not read the leader's
predicted orderings. What it supplied was evidence that a frame the loop had
\emph{closed} was still open. Internal notes had written Task~3 off --- ``nothing to
learn per slot'' --- and the decomposition showed a $0.0661$ C-index gap there, which
is not a statement about our tie-break but about ordering information we were not
using at all. That prompted the re-examination that found PI-RADS. The distinction
matters for anyone repeating this: competitor metrics were useful as a
\emph{falsifier of our own beliefs about saturation}, and would have been actively
harmful used as a target.

\paragraph{A hazard, reported rather than used.}
The same dumps carry a defect. Judge feedback is published verbatim in a
\code{reason} field, and because the judge quotes the entrant's rationale back, those
strings contain clinical quantities from the validation cases: $104$ of $109$ quote
at least one PI-RADS, ISUP, PSA or volume value. Validation inputs are therefore
partially reconstructable from a public leaderboard artefact, and an entry willing to
do so could fit on the validation cohort directly. We did not, and we reported it
alongside the \code{cost\_aware\_tool\_score} property. Both findings have the same
shape --- an agent reading the platform closely enough to find the exploit is also
the thing that can report it --- and we think an agent-run entry has a
\emph{heightened} obligation here, because the marginal cost of noticing such things
is what dropped.

\paragraph{It refuted its own measurements.}
The loop's most valuable single act was catching a bug in its own harness. Comparing
survival models for Task~3, an $L_2$ Cox fitted on CAPRA-S alone scored $0.7249$
against raw CAPRA-S's $0.7372$ --- impossible, since a fitted monotone re-expression
of one variable cannot reorder it. The cause was that pooled out-of-fold
concordance is invalid for refitted models: each fold's coefficients carry their own
scale, so a cross-fold pair compares two rulers. Restricting to within-fold pairs and
pooling numerator and denominator brought the two into exact agreement at $0.7373$
and changed every fitted-model number in the comparison. A second instance: the
shipped map is $\max(1.0, 120 - s\,\cdot\,\text{risk})$, and at the inherited
$s = 8$ the new PI-RADS term pushed nine of 75 cases onto the clamp, where they tie
and the C-index banks $0.5$ on each --- costing $0.0107$, invisible in any aggregate.
Neither was found by a test. Both were found by treating an implausible number as a
defect rather than as noise, which is the discipline Sect.~\ref{sec:human} shows the
loop lacked about \emph{favourable} numbers. The asymmetry is instructive on its own:
the loop interrogated results that looked wrong and banked results that looked good.

\paragraph{Pre-registration became free.}
Every validation slot in this campaign carries an interval stated in a committed
document \emph{before} the submission --- S2's $+0.046$ to $+0.057$, S3's $+0.004$ to
$+0.013$ with a named floor of exactly $0.0000$ --- together with the one variable
changed and the residual risk. This is ordinary good practice that entries rarely
follow, and the reason is cost: writing it down is a separate task competing with the
work. For an agent the artefact is a byproduct of the working process rather than an
overhead on it. The value showed up immediately: S2 came in \emph{above} its band at
$+0.0705$, and because the band existed, the overshoot was legible as a question
rather than a success, which is what produced the prevalence decomposition and the
honest $+0.024$ projection in Sect.~\ref{sec:results}. A result cannot surprise you
if you never said what you expected.

\paragraph{What this does not show.}
None of the above is a controlled comparison. We ran one entry, with one model, one
human and no control arm, and each claim above is of the form ``this happened, and
its cost profile explains why it is rare'' --- not ``an agent will do this.'' The
honest summary is narrower than the section heading: the loop was good at exhaustive,
low-glamour reading; at decomposition when pointed at data; at doubting numbers that
looked wrong; and at producing the paper trail that makes those things checkable. It
was bad, per Sect.~\ref{sec:human}, at doubting numbers that looked right, and at
noticing which frame it was working inside. Those are close to complements, and we
suspect the pairing is not accidental: the same disposition that reads all 1{,}300
lines of the scorer is the one that keeps optimising within the frame the scorer
implies.

% ---------------------------------------------------------------------------
% INTEGRATION NOTE
%
% main.tex is exactly six body pages with references starting on p7; this
% section measures 2.0 pages, so it cannot simply be \input{}.  Three options:
%
% 1. Merge into Sect. 5, retitled "Where the Loop Needed a Human, and What It
%    Bought".  Costs the most cutting but is the strongest structure: the two
%    lists are complements and the closing paragraph above already argues so.
% 2. Keep as its own section and cut for space.  The compressible parts are
%    the pre-registration paragraph (~0.25pp, argued rather than measured)
%    and the forensics prose after the table (~0.3pp).  The table and the
%    78/22 split should not be cut -- that is the thesis measured against an
%    independent system.
% 3. Hold for the extended version, and land only the 78/22 table in Sect. 4,
%    which costs ~0.35pp and carries most of the value.
%
% Note also: if this section ships, Sect. 5's closing sentence calling
% intervention 5 "mildest and most instructive" needs rewording, and the
% still-open question of whether the PI-RADS finding is a sixth intervention
% is partly answered here -- it appears above as the frame-reopening example.
